\scanpage{16}
pour des hom. de corps, par des foncteurs
de ``changement de base'' pour $k\mathbin{\unsure{-|\longrightarrow}}k'$
qui sont d'ailleurs \deleted{pl. fidèles} \inserted{des équiv. de cat.} [mais
pas \deleted{essentiellement} surjectifs]. (Cela provient
du fait que le groupoïde fondamental
\unsure{de'|de} \deleted{\unsure{une}} \inserted{la} multiplicité modulaire pertinente,
ou d'une courbe alg., sur un corps alg.
clos $k$, ne change pas à équiv. près
par l'extension $k\mathbin{\unsure{-|\longrightarrow}}k'$ du corps de base,
avec $k'$ alg. clos) \deleted{par contre il y a \unsure{encore},
si $k'\ne k$, des courbes alg.}
\margin{Mais \unsure{cela ne signifie bien sûr} \inserted{\unsure{pas}} que $\forall$ courbe alg. sur $k'$ soit isom. sur $k'$ à une courbe provenant de $k$ !}

En particulier, le groupe des automorphis-
mes de $k$ opère sur la catégorie précédente,
que j'ai envie d'appeler la \uline{catégorie de
Teichmüller-Galois} associée à $k$. Les
\deleted{objets} \inserted{classes d'iso.} d'objets connexes de cette
catégorie sont en corr. $1$-$1$ avec les types
$(g,\nu)$ des courbes anabéliennes; les groupes
d'automorphismes d'un objet sont essentielle-
ment (en car. $0$ du moins) des complétés
profinis des groupes de Teichmüller st.-ext.
\margin{15}
